\chapter*{Preface}
\addcontentsline{toc}{chapter}{\protect\numberline{}Preface}

This is a book about mathematics. It is also a book about why you should be gentle with the people around you. The connection between those two things is not obvious, and the book takes its whole length to draw it.

The mathematics is algorithmic information theory: results about probability, description, and computation, together with a universal construction built from fair random program bits. The book follows them through one inversion of intuition, then marks the point where theorem ends and interpretation begins. The arguments are not difficult. The formalism is there for those who want it and can be skimmed by those who do not.

The inversion is this. Enumerating every program is simple; identifying the description that matters to a situated observer is hard. The index by which you locate yourself, the system you call the world, even the self that does the locating: each requires a cut. The universal semimeasure weights candidate programs, not realities or observer occurrences. The book's one postulate says no decomposition is view-independently privileged. Whether the describable totality exists as one object, and how a first-person location would be weighted within it, remain open.

This is the third of three books built in dialogue. \textit{Worldlines} developed a block-universe reading of relativity. \textit{Multitudes} developed an Everettian reading of quantum mechanics and its Born weights. This book turns to program descriptions and asks a different question: what belongs to the mathematics, what belongs to a chosen representation, and what belongs to us?

The formalism supplies no moral preference. We do.

No program weight tells you whom to love or how to act. Care is made here, within causal reach, between finite people who will not meet in exactly this way again. The measure does not provide it. We provide it.

Hello, fellow sufferer. Stay with me.
